\documentclass[10pt,a4paper]{article}
\usepackage[margin=1in]{geometry}
\usepackage{booktabs}
\usepackage{amsmath,amssymb}
\usepackage{enumitem}
\usepackage{graphicx}
\usepackage{xurl}
\usepackage[colorlinks=true,linkcolor=blue,urlcolor=blue]{hyperref}

\title{Experimental Section}
\author{}
\date{\today}

\begin{document}
\maketitle

\section{Experiments}
\label{sec:experiments}

\subsection{Overview}

This section evaluates the proposed transfer architecture under the current paper protocol and answers three questions:
\begin{enumerate}[leftmargin=1.4em]
  \item whether topology-conditioned transfer outperforms simpler shared prediction baselines;
  \item which components in TCDP contribute stable gains under a fixed protocol;
  \item whether observations are consistent across two topology-expert capacities.
\end{enumerate}

\subsection{Experimental Protocol}

All experiments use the same target-domain split protocol:
\begin{itemize}[leftmargin=1.2em]
  \item split mode: \texttt{table\_group},
  \item ratio: train:val$=1:5$,
  \item no independent test split at this stage.
\end{itemize}
The dataset is
\path{paper_materials/data/dataset_table_group_train_val_no_test.pkl}.
Therefore, this section focuses on model selection and mechanism validation under train/val evidence.

\subsection{Implementation Details}

Unless explicitly ablated, the training setup is fixed:
\begin{itemize}[leftmargin=1.2em]
  \item script: \path{src/train_balanced_sampling_sep_mlp_shared_calib_step5a_feat_opt.py},
  \item frozen HGAT with cached structural embeddings (\texttt{freeze\_hgat}),
  \item AdamW optimizer with plateau LR scheduler,
  \item batch size 2048, learning rate $2\times10^{-3}$, weight decay $5\times10^{-5}$,
  \item \texttt{mlp\_dropout=0.2}, \texttt{z\_noise\_std=0.01},
  \item early stop: patience 12, min delta $8\times10^{-4}$,
  \item validation every 5 epochs.
\end{itemize}
All reported tables aggregate five seeds:
\[
\{9294,9295,9296,9297,9298\}.
\]

\subsection{Metrics and Aggregation}

For each seed, we parse the \texttt{[Best]} checkpoint selected by validation loss and record
\texttt{Train R\textsuperscript{2}@Best}, \texttt{Val R\textsuperscript{2}}, and \texttt{Best Epoch}.
The primary metric is
\begin{equation}
R^2 = 1-\frac{\sum_{i}(y_i-\hat{y}_i)^2}{\sum_{i}(y_i-\bar{y})^2},
\end{equation}
Across $K=5$ seeds, each metric $m$ is summarized as
\begin{equation}
\mu(m)=\frac{1}{K}\sum_{k=1}^K m_k,\qquad
\sigma(m)=\sqrt{\frac{1}{K-1}\sum_{k=1}^K (m_k-\mu(m))^2}.
\end{equation}

\subsection{Core Model Comparison}

Table~\ref{tab:core_compare_en} reports the core architecture comparison.

\begin{table}[htbp]
\centering
\small
\caption{Core comparison under \texttt{table\_group}, train:val$=1:5$ (no test split).}
\label{tab:core_compare_en}
\begin{tabular}{lccc}
\toprule
Method & NSeed & Train $R^2$ & Val $R^2$ \\
\midrule
Baseline(shared linear) & 5 & $0.9024\pm0.0086$ & $0.8726\pm0.0067$ \\
TCDP(16,96) & 5 & $0.9586\pm0.0085$ & $0.9065\pm0.0168$ \\
TCDP(32,64) & 5 & $0.9582\pm0.0216$ & $\mathbf{0.9234\pm0.0175}$ \\
\bottomrule
\end{tabular}
\end{table}

Compared with the shared-linear baseline, TCDP(32,64) improves Val $R^2$ by +0.0508
Compared with TCDP(16,96), TCDP(32,64) further improves Val $R^2$ by +0.0169.

\subsection{One-Factor Ablation on TCDP(32,64)}

Under a one-factor protocol, we modify exactly one component while keeping all other settings fixed.
Results are listed in Table~\ref{tab:t32_ablation_en}.

\begin{table}[htbp]
\centering
\scriptsize
\setlength{\tabcolsep}{3.5pt}
\caption{One-factor ablation on TCDP(32,64) (5 seeds each).}
\label{tab:t32_ablation_en}
\resizebox{\columnwidth}{!}{
\begin{tabular}{lccc}
\toprule
Variant & NSeed & Train $R^2$ & Val $R^2$ \\
\midrule
A1 (MB-enrich on) & 5 & $0.9682\pm0.0087$ & $\mathbf{0.9298\pm0.0073}$ \\
A2 (EnrichOff-auto) & 5 & $0.9682\pm0.0087$ & $\mathbf{0.9298\pm0.0073}$ \\
A3 (ParasiticReplace) & 5 & $0.9694\pm0.0080$ & $0.9294\pm0.0084$ \\
A4 (DualMean) & 5 & $0.9582\pm0.0216$ & $0.9234\pm0.0175$ \\
A5 (MB + TopoOff) & 5 & $0.9006\pm0.0068$ & $0.8720\pm0.0034$ \\
\bottomrule
\end{tabular}
}
\vspace{2pt}
\raggedright
\footnotesize
A1: T32x64 + ModeBase(enrich\_on).\\
A2: T32x64 + EnrichOff(auto).\\
A3: T32x64 + ParasiticReplace.\\
A4: T32x64 + DualMean.\\
A5: T32x64 + ModeBase+TopoOff.
\end{table}

Three observations are consistent:
\begin{enumerate}[leftmargin=1.4em]
  \item Turning topology expert off causes the largest collapse
  ($\Delta R^2 \approx -0.0578$), validating topology-conditioned residual correction as the dominant gain source.
  \item A3 is neutral-to-slightly negative versus the best ModeBase path.
  \item A1 and A2 are numerically identical in this implementation path and should be treated as an equivalence-control pair.
\end{enumerate}

\subsection{Supplementary Ablation on TCDP(16,96)}

To verify trend stability across capacity settings, Table~\ref{tab:t16_ablation_en}
reports one-factor ablations for TCDP(16,96).

\begin{table}[htbp]
\centering
\scriptsize
\setlength{\tabcolsep}{3.5pt}
\caption{Supplementary one-factor ablation on TCDP(16,96) (5 seeds each).}
\label{tab:t16_ablation_en}
\resizebox{\columnwidth}{!}{
\begin{tabular}{lccc}
\toprule
Variant & NSeed & Train $R^2$ & Val $R^2$ \\
\midrule
B1 (MB-enrich on) & 5 & $0.9548\pm0.0122$ & $\mathbf{0.9160\pm0.0139}$ \\
B2 (EnrichOff-auto) & 5 & $0.9548\pm0.0122$ & $\mathbf{0.9160\pm0.0139}$ \\
B3 (ParasiticReplace) & 5 & $0.9562\pm0.0118$ & $0.9151\pm0.0113$ \\
B4 (DualMean) & 5 & $0.9586\pm0.0085$ & $0.9065\pm0.0168$ \\
B5 (MB + TopologyOff) & 5 & $0.9006\pm0.0068$ & $0.8720\pm0.0034$ \\
\bottomrule
\end{tabular}
}
\vspace{2pt}
\raggedright
\footnotesize
B1: T16x96 + mode\_base(enrich\_on).\\
B2: T16x96 + enrich\_off(auto).\\
B3: T16x96 + parasitic\_replace.\\
B4: T16x96 + dual\_readout(mean).\\
B5: T16x96 + mode\_base+topology\_off.
\end{table}

The same pattern remains: topology-off gives the strongest degradation,
while parasitic replacement and dual-readout do not provide stable improvements.

\end{document}

